\chapter{Dictionaries and Sets}

\textbf{What this chapter gives you:} the two containers that find things by \emph{value} instead of by
position. A \code{dict} maps keys to values and is the shape of nearly all JSON, API requests, model replies,
and configuration. A \code{set} is a bag of unique items with fast membership tests and set algebra. After
this chapter you can read any nested API response and pull out one number safely.

%RELATED: 6, 7, 8, 10, 11, 12

% ------------------------------------------------------------
\section{Dictionaries}

A dictionary stores \emph{key $\to$ value} pairs. You ask for a value by its key, not by a position.
Keys are unique. Since Python 3.7 a dict also remembers the order in which keys were added.

%MD<
% ```text
%   list: find by POSITION            dict: find by KEY
%
%   ["haiku", "sonnet"]               {"fast": "haiku", "smart": "sonnet"}
%     0        1                        "fast"   -> "haiku"
%     |        |                        "smart"  -> "sonnet"
%   m[0]     m[1]                     m["fast"]   m["smart"]
% ```
%MD>
\begin{center}
\begin{tikzpicture}[k/.style={draw=cgold,fill=cgold!15,minimum width=14mm,minimum height=7mm,font=\ttfamily\small},
  v/.style={draw=cblue,fill=cblue!8,minimum width=18mm,minimum height=7mm,font=\ttfamily\small},
  lab/.style={font=\sffamily\scriptsize,text=cgray},>=Stealth]
  \node[lab] at (1.9,2.2) {\textbf{list}: look up by position};
  \node[k] (i0) at (0.2,1.3) {0};  \node[v] (v0) at (2.4,1.3) {"haiku"};
  \node[k] (i1) at (0.2,0.4) {1};  \node[v] (v1) at (2.4,0.4) {"sonnet"};
  \draw[->,cgray] (i0) -- (v0);   \draw[->,cgray] (i1) -- (v1);
  \node[lab] at (1.3,-0.4) {\code{m[1]} gives \code{"sonnet"}};
  \node[lab] at (8.4,2.2) {\textbf{dict}: look up by key};
  \node[k,minimum width=18mm] (k0) at (6.6,1.3) {"fast"};  \node[v] (w0) at (9.4,1.3) {"haiku"};
  \node[k,minimum width=18mm] (k1) at (6.6,0.4) {"smart"}; \node[v] (w1) at (9.4,0.4) {"sonnet"};
  \draw[->,cgray] (k0) -- (w0);   \draw[->,cgray] (k1) -- (w1);
  \node[lab] at (8.0,-0.4) {\code{m["smart"]} gives \code{"sonnet"}};
\end{tikzpicture}
\end{center}

\subsection{Create with \{\} and dict()}

\begin{pycode}
>>> usage = {"input_tokens": 812, "output_tokens": 143}
>>> usage["input_tokens"]
812
>>> same = dict(input_tokens=812, output_tokens=143)
>>> same == usage
True
>>> type(usage), len(usage), {}
(<class 'dict'>, 2, {})
\end{pycode}

\code{dict(name=value)} is handy, but the names must be legal variable names (no spaces, no hyphens, no
reserved words). For keys like \code{"max-tokens"}, use the brace form. A trailing comma after the last
pair is allowed.

\subsection{Convert with dict()}

\code{dict()} also accepts a sequence of two-item sequences. \code{zip()} (Chapter 8) is the usual way to
make them.

\begin{pycode}
>>> dict([("a", 1), ("b", 2)])
{'a': 1, 'b': 2}
>>> dict(["ab", "cd"])
{'a': 'b', 'c': 'd'}
>>> dict(zip(["x", "y"], [10, 20]))
{'x': 10, 'y': 20}
>>> dict([("a", 1, 2)])
Traceback (most recent call last):
  ...
ValueError: dictionary update sequence element #0 has length 3; 2 is required
\end{pycode}

\subsection{Add or change an item with [key]}

Assigning to a key adds it when it is new and replaces the value when it exists. There is no
out-of-range error as with lists. Duplicate keys in a literal keep the last value.

\begin{pycode}
>>> cfg = {"model": "haiku"}
>>> cfg["max_tokens"] = 1000
>>> cfg["model"] = "sonnet"
>>> cfg
{'model': 'sonnet', 'max_tokens': 1000}
>>> {"a": 1, "a": 2}
{'a': 2}
>>> list(cfg)
['model', 'max_tokens']
\end{pycode}

\begin{typecheck}[: what can be a key]
A key must be \emph{hashable}, which in practice means immutable: \code{str}, \code{int}, \code{float},
\code{bool}, \code{None}, \code{bytes}, \code{frozenset}, and tuples whose items are hashable. A list, dict,
or set cannot be a key. Values can be anything.
\end{typecheck}

\begin{pycode}
>>> {(1, 2): "tuple key works"}
{(1, 2): 'tuple key works'}
>>> {[1, 2]: "no"}
Traceback (most recent call last):
  ...
TypeError: unhashable type: 'list'
\end{pycode}

\begin{trap}[: keys that are equal count as the same key]
\code{1}, \code{1.0}, and \code{True} are equal, so they collide. The first key object stays, and the last
value wins.
\end{trap}

\begin{pycode}
>>> {1: "a", True: "b", 1.0: "c"}
{1: 'c'}
\end{pycode}

\subsection{Get an item by [key] or with get()}

\begin{center}
\begin{tabularx}{\linewidth}{@{}l l X@{}}
\toprule
\textbf{Expression} & \textbf{Key present} & \textbf{Key missing} \\
\midrule
\code{d[k]}           & the value & raises \code{KeyError} \\
\code{d.get(k)}       & the value & returns \code{None} \\
\code{d.get(k, dflt)} & the value & returns \code{dflt} \\
\code{k in d}         & \code{True} & \code{False} \\
\code{d.pop(k)}       & value, and removes the pair & raises \code{KeyError} \\
\code{d.pop(k, dflt)} & value, and removes the pair & returns \code{dflt} \\
\code{d.setdefault(k, dflt)} & the value & stores \code{dflt} under \code{k}, then returns it \\
\bottomrule
\end{tabularx}
\end{center}

\begin{pycode}
>>> resp = {"model": "claude-haiku-5-5", "stop_reason": "end_turn"}
>>> resp["model"]
'claude-haiku-5-5'
>>> resp["usage"]
Traceback (most recent call last):
  ...
KeyError: 'usage'
>>> resp.get("usage") is None
True
>>> resp.get("usage", {})
{}
>>> "usage" in resp
False
\end{pycode}

\begin{trap}[: \texttt{d.key} is not how dicts work]
A dict is read with \emph{brackets} and a key: \code{d["model"]}. Dot access (\code{d.model}) looks for an
\emph{attribute} and fails. SDK response objects are the opposite: they use dots
(\code{response.usage}), while data parsed from JSON is a dict and uses brackets.
\end{trap}

\begin{pycode}
>>> resp = {"model": "claude-haiku-5-5"}
>>> resp.model
Traceback (most recent call last):
  ...
AttributeError: 'dict' object has no attribute 'model'
\end{pycode}

\begin{center}
\begin{tabularx}{\linewidth}{@{}l l X@{}}
\toprule
\textbf{Syntax} & \textbf{Looks up} & \textbf{Fails with} \\
\midrule
\code{d["k"]}  & a dict key         & \code{KeyError} \\
\code{obj.x}   & an attribute       & \code{AttributeError} \\
\code{lst[0]}  & a list position    & \code{IndexError} \\
\code{d.get("k")} & a dict key, safely & never; gives \code{None} \\
\bottomrule
\end{tabularx}
\end{center}

\subsection{Nested dictionaries: reading JSON}

Values can be dicts and lists, so data nests as deeply as you like. Read it with one subscript per level,
left to right. The picture shows how \code{reply["content"][0]["text"]} walks down the tree.

%MD<
% ```mermaid
% flowchart LR
%     reply{{"reply (dict)"}} --> id["'id': 'msg_01'"]
%     reply --> usage{{"'usage' (dict)"}}
%     usage --> i1["'input_tokens': 812"]
%     usage --> i2["'output_tokens': 143"]
%     reply -- "['content']" --> content[["'content' (list)"]]
%     content -- "[0]" --> item{{"item 0 (dict)"}}
%     item --> ty["'type': 'text'"]
%     item -- "['text']" --> tx["'text': 'Hello'"]
% ```
%MD>
\begin{center}
\begin{tikzpicture}[dict/.style={draw=cblue,fill=cblue!10,rounded corners=2pt,font=\ttfamily\scriptsize,inner sep=3pt},
  lst/.style={draw=cgold,fill=cgold!15,rounded corners=2pt,font=\ttfamily\scriptsize,inner sep=3pt},
  leaf/.style={draw=cgreen,fill=cgreen!10,rounded corners=2pt,font=\ttfamily\scriptsize,inner sep=3pt},
  ed/.style={cgray},hot/.style={thick,cred},
  el/.style={font=\ttfamily\tiny,text=cred,fill=white,inner sep=1pt},>=Stealth]
  \node[dict] (r)  at (0,0) {reply (dict)};
  \node[leaf] (id) at (4.2,1.9) {"id": "msg\_01"};
  \node[dict] (us) at (4.2,0.5) {"usage" (dict)};
  \node[lst]  (co) at (4.2,-1.4) {"content" (list)};
  \node[leaf] (i1) at (8.0,1.0) {"input\_tokens": 812};
  \node[leaf] (i2) at (8.0,0.0) {"output\_tokens": 143};
  \node[dict] (d0) at (8.0,-1.4) {item 0 (dict)};
  \node[leaf] (ty) at (11.7,-0.9) {"type": "text"};
  \node[leaf] (tx) at (11.7,-1.9) {"text": "Hello"};
  \draw[->,ed] (r) -- (id);
  \draw[->,ed] (r) -- (us);
  \draw[->,hot] (r) -- node[el,midway]{["content"]} (co);
  \draw[->,ed] (us) -- (i1);
  \draw[->,ed] (us) -- (i2);
  \draw[->,hot] (co) -- node[el,midway]{[0]} (d0);
  \draw[->,ed] (d0) -- (ty);
  \draw[->,hot] (d0) -- node[el,midway]{["text"]} (tx);
\end{tikzpicture}
\end{center}

\begin{pycode}
>>> reply = {
...     "id": "msg_01",
...     "usage": {"input_tokens": 812, "output_tokens": 143},
...     "content": [{"type": "text", "text": "Hello"}],
... }
>>> reply["usage"]["input_tokens"]
812
>>> reply["content"][0]["text"]
'Hello'
>>> reply["usage"]["total"]
Traceback (most recent call last):
  ...
KeyError: 'total'
>>> reply.get("error", {}).get("message", "no error")
'no error'
\end{pycode}

\begin{boss}[: missing keys deep inside nested data]
Real replies are not always complete. A key can be absent, or present with the value \code{None}, or a list
can be empty. \code{reply["a"]["b"]["c"]} fails at the first gap with a \code{KeyError} or \code{TypeError}
that does not say which level was wrong.\\[2pt]
\textbf{Strategy:} decide per level. Use \code{[ ]} when absence is a bug (you want the loud error). Use
\code{.get(key, \{\})} to step through levels that may be missing, then \code{.get(key, default)} at the last
one. Check lists with \code{if items:} before \code{items[0]}. Remember \code{.get(k, d)} returns \code{d}
only when the key is \emph{missing}, not when its value is \code{None}.
\end{boss}

\begin{pycode}
>>> reply = {"usage": None, "content": []}
>>> reply.get("usage", {}).get("input_tokens", 0)
Traceback (most recent call last):
  ...
AttributeError: 'NoneType' object has no attribute 'get'
>>> (reply.get("usage") or {}).get("input_tokens", 0)
0
>>> first = reply["content"][0] if reply["content"] else None
>>> first is None
True
\end{pycode}

\subsection{Iterate with for and in}

Looping over a dict gives its \emph{keys}, in insertion order. \code{.keys()}, \code{.values()}, and
\code{.items()} return \emph{views}: light objects that follow the dict and can be looped over or
turned into lists.

\begin{pycode}
>>> tokens = {"prompt": 812, "completion": 143}
>>> for name in tokens:
...     print(name)
prompt
completion
>>> list(tokens.keys()), list(tokens.values())
(['prompt', 'completion'], [812, 143])
>>> list(tokens.items())
[('prompt', 812), ('completion', 143)]
>>> for name, n in tokens.items():
...     print(name, "->", n)
prompt -> 812
completion -> 143
>>> sum(tokens.values())
955
\end{pycode}

\begin{typecheck}[: what the loop variable is]
\code{for k in d} and \code{for k in d.keys()}: the key. \code{for v in d.values()}: the value.
\code{for item in d.items()}: a \code{(key, value)} \emph{tuple}, so unpack it in the \code{for} line
(Chapter 8) with \code{for k, v in d.items()}.
\end{typecheck}

\begin{trap}[: do not add or remove keys while looping]
Changing the number of keys during a loop raises \code{RuntimeError}. Loop over a copy of the keys
(\code{list(d)}) when you must change the dict, or build a new dict.
\end{trap}

\begin{pycode}
>>> d = {"a": 1, "b": 2}
>>> for k in d:
...     d[k + "x"] = 0
Traceback (most recent call last):
  ...
RuntimeError: dictionary changed size during iteration
>>> for k in list(d):
...     d[k + "x"] = 0
>>> sorted(d)
['a', 'ax', 'axx', 'b', 'bx']
\end{pycode}

\subsection{Combine dictionaries}

When two dicts share a key, the one applied \emph{last} wins. There are three ways to merge.

\begin{center}
\begin{tabularx}{\linewidth}{@{}l l X@{}}
\toprule
\textbf{Way} & \textbf{Result} & \textbf{Changes the left dict?} \\
\midrule
\code{a | b}          & new dict & no \\
\code{\{**a, **b\}}   & new dict & no \\
\code{a.update(b)}    & \code{None} & yes, in place \\
\code{a |= b}         & (assignment) & yes, in place \\
\bottomrule
\end{tabularx}
\end{center}

\begin{pycode}
>>> base = {"model": "haiku", "temperature": 0.7}
>>> override = {"temperature": 0.0, "max_tokens": 500}
>>> base | override
{'model': 'haiku', 'temperature': 0.0, 'max_tokens': 500}
>>> {**base, **override} == base | override
True
>>> base
{'model': 'haiku', 'temperature': 0.7}
>>> base.update(override)
>>> base["max_tokens"]
500
\end{pycode}

This ``defaults, then overrides'' pattern is how you combine default request settings with per-call
settings. All three ways are shallow: nested values are shared, not copied.

\subsection{Delete items}

\begin{pycode}
>>> d = {"a": 1, "b": 2, "c": 3}
>>> del d["a"]
>>> d.pop("b")
2
>>> d.pop("zzz", None) is None
True
>>> d.popitem()
('c', 3)
>>> d
{}
>>> del d["x"]
Traceback (most recent call last):
  ...
KeyError: 'x'
\end{pycode}

\code{popitem()} removes and returns the most recently added pair. \code{clear()} empties the dict in place;
assigning \code{d = \{\}} makes a new empty dict and leaves other names pointing at the old one.

\subsection{Test for a key with in}

\code{in} on a dict checks \emph{keys}, never values. Test values with \code{v in d.values()}, which scans.

\begin{pycode}
>>> limits = {"input": 100, "output": 20}
>>> "input" in limits, 100 in limits, 100 in limits.values()
(True, False, True)
\end{pycode}

\subsection{Assign, copy, deepcopy}

Dicts are mutable, so everything from Chapter 8 applies. \code{b = a} is a second name for one dict.
\code{a.copy()} is shallow: new outer dict, shared inner values. \code{copy.deepcopy(a)} copies all the way
down.

\begin{pycode}
>>> import copy
>>> a = {"model": "haiku", "stops": ["END"]}
>>> b = a
>>> c = a.copy()
>>> e = copy.deepcopy(a)
>>> a["stops"].append("STOP")
>>> c["stops"], e["stops"]
(['END', 'STOP'], ['END'])
>>> b is a, c is a, c["stops"] is a["stops"]
(True, False, True)
\end{pycode}

\subsection{Compare dictionaries}

\code{==} and \code{!=} compare keys and values and ignore order. The ordering operators
(\code{<}, \code{>}) are not defined for dicts.

\begin{pycode}
>>> {1: 1, 2: 2, 3: 3} == {3: 3, 1: 1, 2: 2}
True
>>> {1: [1, 2]} == {1: [1, 1]}
False
>>> {1: 1} < {2: 2}
Traceback (most recent call last):
  ...
TypeError: '<' not supported between instances of 'dict' and 'dict'
\end{pycode}

\subsection{Dictionary comprehensions}

The form is \code{\{key\_expr: value\_expr for item in iterable if condition\}}. Use it to build or reshape
dicts in one expression.

\begin{pycode}
>>> words = ["haiku", "opus", "sonnet"]
>>> {w: len(w) for w in words}
{'haiku': 5, 'opus': 4, 'sonnet': 6}
>>> {w: len(w) for w in words if len(w) > 4}
{'haiku': 5, 'sonnet': 6}
>>> {v: k for k, v in {"a": 1, "b": 2}.items()}
{1: 'a', 2: 'b'}
\end{pycode}

Swapping keys and values drops data when values repeat, and fails when a value is not hashable.

\subsection{Counting and grouping patterns}

Two patterns appear constantly. Counting: \code{d.get(key, 0) + 1}. Grouping: \code{d.setdefault(key, [])}
then \code{.append}. For counting alone, \code{collections.Counter} does it in one call (Chapter 12).

\begin{pycode}
>>> counts = {}
>>> for ch in "letters":
...     counts[ch] = counts.get(ch, 0) + 1
>>> counts
{'l': 1, 'e': 2, 't': 2, 'r': 1, 's': 1}
>>> by_label = {}
>>> for doc, label in [("d1", "spam"), ("d2", "ok"), ("d3", "spam")]:
...     by_label.setdefault(label, []).append(doc)
>>> by_label
{'spam': ['d1', 'd3'], 'ok': ['d2']}
>>> from collections import Counter
>>> Counter("letters").most_common(1)
[('e', 2)]
\end{pycode}

% ------------------------------------------------------------
\section{Sets}

A set is a dict with the values thrown away: a bag of \emph{unique, hashable} items with no order and no
positions. Use one when you only care whether something is in there, or when you need to remove duplicates
or compare groups.

\subsection{Create with set() or \{\}}

\begin{pycode}
>>> type({}), type(set())
(<class 'dict'>, <class 'set'>)
>>> seen = {3, 1, 3, 2}
>>> seen
{1, 2, 3}
>>> sorted(set("letters"))
['e', 'l', 'r', 's', 't']
>>> set((1,)), set([0, 2, 4]), set()
({1}, {0, 2, 4}, set())
\end{pycode}

\begin{trap}[: \texttt{\{\}} is an empty dict, not an empty set]
Empty braces are a \emph{dict}. Write \code{set()} for an empty set. Also \code{set()} takes \emph{one}
iterable, not several values, and a bare number is not iterable.
\end{trap}

\begin{pycode}
>>> set(0, 2, 4)
Traceback (most recent call last):
  ...
TypeError: set expected at most 1 argument, got 3
>>> set(5)
Traceback (most recent call last):
  ...
TypeError: 'int' object is not iterable
\end{pycode}

Because a set is unordered, you cannot index it, and the order of a printed set of strings can change
between runs. Use \code{sorted(s)} when you need a stable order. Items must be hashable.

\begin{pycode}
>>> {1, 2}[0]
Traceback (most recent call last):
  ...
TypeError: 'set' object is not subscriptable
>>> {[1, 2]}
Traceback (most recent call last):
  ...
TypeError: unhashable type: 'list'
\end{pycode}

\subsection{Add, remove, length, membership}

\begin{pycode}
>>> s = {1, 2, 3}
>>> s.add(4)
>>> s.add(2)
>>> s.remove(3)
>>> s.discard(99)
>>> s
{1, 2, 4}
>>> len(s), 2 in s, 3 in s
(3, True, False)
>>> s.remove(99)
Traceback (most recent call last):
  ...
KeyError: 99
\end{pycode}

\code{remove()} raises \code{KeyError} when the item is absent. \code{discard()} stays quiet. Adding an item
that is already there changes nothing. Testing \code{x in s} takes about the same time however big the set
is; on a list it scans every item.

A common use is removing duplicates. A set loses the order; a dict keeps it:

\begin{pycode}
>>> docs = ["b", "a", "b", "c", "a"]
>>> len(set(docs))
3
>>> list(dict.fromkeys(docs))
['b', 'a', 'c']
\end{pycode}

\subsection{Set operations}

Set algebra answers ``what do these two groups share?'' and ``what is missing?''. Each operator has a
method twin. The methods accept any iterable; the operators need two sets.

%MD<
% ```text
%   a = {1, 2}     b = {2, 3}
%
%   a | b  union          a & b  intersection      a - b  difference       a ^ b  symmetric diff.
%   ( 1 [2] 3 )           ( 1  [2]  3 )            ( [1] 2   3 )           ( [1]  2  [3] )
%   all items             items in both            in a, not in b          in exactly one
%   {1, 2, 3}             {2}                      {1}                     {1, 3}
% ```
%MD>
\begin{center}
\begin{tikzpicture}[lab/.style={font=\ttfamily\small,align=center},it/.style={font=\ttfamily\scriptsize}]
  % panel 1: union
  \begin{scope}[shift={(0,0)}]
    \fill[cblue,opacity=0.35] (-0.4,0) circle (0.7);  \fill[cblue,opacity=0.35] (0.4,0) circle (0.7);
    \draw[cgray] (-0.4,0) circle (0.7);  \draw[cgray] (0.4,0) circle (0.7);
    \node[it] at (-0.75,0) {1}; \node[it] at (0,0) {2}; \node[it] at (0.75,0) {3};
    \node[lab] at (0,-1.3) {a | b\\ \footnotesize\sffamily union};
  \end{scope}
  % panel 2: intersection
  \begin{scope}[shift={(4.0,0)}]
    \begin{scope}\clip (-0.4,0) circle (0.7); \fill[cgreen,opacity=0.45] (0.4,0) circle (0.7);\end{scope}
    \draw[cgray] (-0.4,0) circle (0.7);  \draw[cgray] (0.4,0) circle (0.7);
    \node[it] at (-0.75,0) {1}; \node[it] at (0,0) {2}; \node[it] at (0.75,0) {3};
    \node[lab] at (0,-1.3) {a \& b\\ \footnotesize\sffamily intersection};
  \end{scope}
  % panel 3: difference
  \begin{scope}[shift={(8.0,0)}]
    \fill[cgold,opacity=0.45] (-0.4,0) circle (0.7);  \fill[white] (0.4,0) circle (0.7);
    \draw[cgray] (-0.4,0) circle (0.7);  \draw[cgray] (0.4,0) circle (0.7);
    \node[it] at (-0.75,0) {1}; \node[it] at (0,0) {2}; \node[it] at (0.75,0) {3};
    \node[lab] at (0,-1.3) {a - b\\ \footnotesize\sffamily difference};
  \end{scope}
  % panel 4: symmetric difference
  \begin{scope}[shift={(12.0,0)}]
    \fill[cpurple,opacity=0.4] (-0.4,0) circle (0.7);  \fill[cpurple,opacity=0.4] (0.4,0) circle (0.7);
    \begin{scope}\clip (-0.4,0) circle (0.7); \fill[white] (0.4,0) circle (0.7);\end{scope}
    \draw[cgray] (-0.4,0) circle (0.7);  \draw[cgray] (0.4,0) circle (0.7);
    \node[it] at (-0.75,0) {1}; \node[it] at (0,0) {2}; \node[it] at (0.75,0) {3};
    \node[lab] at (0,-1.3) {a \textasciicircum{} b\\ \footnotesize\sffamily in exactly one};
  \end{scope}
  \node[font=\sffamily\scriptsize,text=cgray] at (6.0,1.15) {with \code{a = \{1, 2\}} (left circle) and \code{b = \{2, 3\}} (right circle); the shaded part is the result};
\end{tikzpicture}
\end{center}

\begin{center}
\begin{tabularx}{\linewidth}{@{}l l X@{}}
\toprule
\textbf{Operator} & \textbf{Method} & \textbf{Result is} \\
\midrule
\code{a | b}  & \code{a.union(b)}                & items in either \\
\code{a \& b} & \code{a.intersection(b)}         & items in both \\
\code{a - b}  & \code{a.difference(b)}           & items in \code{a} but not in \code{b} \\
\code{a \textasciicircum{} b} & \code{a.symmetric\_difference(b)} & items in exactly one \\
\code{a <= b} & \code{a.issubset(b)}             & \code{bool}: every item of \code{a} is in \code{b} \\
\code{a < b}  & --                               & \code{bool}: proper subset (and not equal) \\
\code{a >= b} & \code{a.issuperset(b)}           & \code{bool}: \code{a} contains all of \code{b} \\
--            & \code{a.isdisjoint(b)}           & \code{bool}: nothing in common \\
\bottomrule
\end{tabularx}
\end{center}

\begin{pycode}
>>> a, b = {1, 2}, {2, 3}
>>> a | b, a & b, a - b, a ^ b
({1, 2, 3}, {2}, {1}, {1, 3})
>>> allowed = {"search", "calc", "email"}
>>> requested = {"calc", "shell"}
>>> sorted(requested & allowed), sorted(requested - allowed)
(['calc'], ['shell'])
>>> requested <= allowed, {1} <= {1, 2}, {1, 2} < {1, 2}
(False, True, False)
>>> a.union([5, 6]) == {1, 2, 5, 6}
True
>>> a | [5, 6]
Traceback (most recent call last):
  ...
TypeError: unsupported operand type(s) for |: 'set' and 'list'
\end{pycode}

The result of an operator is a new set, and an empty result is falsy, so
\code{if requested - allowed:} reads as ``is anything missing?''.

\subsection{Set comprehensions and frozenset}

\begin{pycode}
>>> {n % 3 for n in range(10)}
{0, 1, 2}
>>> sorted({w[0] for w in ["apple", "avocado", "banana"]})
['a', 'b']
>>> fs = frozenset([3, 2, 1])
>>> fs
frozenset({1, 2, 3})
>>> fs.add(4)
Traceback (most recent call last):
  ...
AttributeError: 'frozenset' object has no attribute 'add'
>>> {frozenset({1, 2}): "usable as a dict key"}
{frozenset({1, 2}): 'usable as a dict key'}
\end{pycode}

A \code{frozenset} is an immutable set. Because it is hashable it can be a dict key or a member of another
set, which an ordinary set cannot.

\begin{typecheck}[: the four comprehensions and the odd one out]
\begin{center}
\begin{tabularx}{\linewidth}{@{}l l l X@{}}
\toprule
\textbf{Syntax} & \textbf{Builds a} & \textbf{Duplicates} & \textbf{Note} \\
\midrule
\code{[x for ...]}         & \code{list}      & kept            & ordered \\
\code{\{x for ...\}}       & \code{set}       & dropped         & unordered \\
\code{\{k: v for ...\}}    & \code{dict}      & later key wins  & the colon makes it a dict \\
\code{(x for ...)}         & generator        & kept            & lazy, one pass; \emph{not} a tuple \\
\bottomrule
\end{tabularx}
\end{center}
\end{typecheck}

\subsection{Which container?}

\begin{center}
\begin{tabularx}{\linewidth}{@{}l c c c c@{}}
\toprule
 & \textbf{tuple} & \textbf{list} & \textbf{dict} & \textbf{set} \\
\midrule
Ordered?          & \cmark & \cmark & \cmark (by insertion) & \xmark \\
Duplicates?       & \cmark & \cmark & keys: \xmark & \xmark \\
Mutable?          & \xmark & \cmark & \cmark & \cmark \\
Find by           & position & position & key & value (\code{in}) \\
Fast \code{x in}? & \xmark & \xmark & \cmark (keys) & \cmark \\
\bottomrule
\end{tabularx}
\end{center}

\begin{gamerule}
Reach for a dict when you have named fields or a lookup table, a set when you ask ``have I seen this?'',
a list when order and repeats matter, and a tuple for a small fixed record.
\end{gamerule}

\begin{practice}
\begin{enumerate}
\item Build a dict \code{prices} mapping \code{"haiku"}, \code{"sonnet"}, \code{"opus"} to three prices. Print
the price of \code{"sonnet"} two ways, one of which cannot raise \code{KeyError}.
\item Given \code{prices}, build the inverse dict (price to name). What goes wrong if two models share a price?
\item What is the difference between \code{d.get("k")} returning \code{None} and \code{d["k"]}?
\item Merge \code{\{"a": 1, "b": 2\}} and \code{\{"b": 9\}} into a new dict without changing either. Which value
does \code{"b"} get?
\item Why does \code{\{\}} not create an empty set, and what do you write instead?
\item Given two sets of tool names, \code{allowed} and \code{requested}, write one expression for the tools
that were requested but not allowed.
\end{enumerate}
\end{practice}

\begin{achievement}{Key Master}
You can now:
\begin{itemize}
\item create, read, update, merge, and delete dict entries, choosing between \code{[ ]}, \code{get}, and
\code{in};
\item walk nested JSON-like data and handle missing keys safely;
\item loop with \code{keys()}, \code{values()}, and \code{items()} without changing the dict during the loop;
\item use sets for uniqueness, membership, and \code{|}, \code{\&}, \code{-}, \code{\textasciicircum{}};
\item explain why a list cannot be a dict key and when to use \code{frozenset}.
\end{itemize}
\end{achievement}
